We revisit the combinatorial loop-equation method for computing planar disc amplitudes in
polynomial multi-matrix models. We develop a procedure which organises the Schwinger-Dyson
constraints by the level of their boundary insertions and, when the resulting system closes,
produces an algebraic equation for the planar resolvent. Using this method we recover the disc
partition functions associated with fixed, mixed, and free boundary conditions in the three-state
Potts model, previously obtained by saddle-point methods. We further show one can compute the disc
partition function for several unsolved two- and three- matrix models. These examples suggest that
loop equations provide a practical route to identifying and solving tractable sectors of
multi-matrix models beyond the standard one- and two-matrix cases.